\documentclass{article}
\usepackage{xfp}
\usepackage{tikzplot}

\begin{document}

% \tikzmath{
%   % 5 点的齐次坐标
%   \A{1,1} = -2; \A{2,1} = 2; \A{3,1} = 1; 
%   \B{1,1} = 1; \B{2,1} = 2; \B{3,1} = 1; 
%   \C{1,1} = 2; \C{2,1} = 0; \C{3,1} = 1; 
%   \D{1,1} = -1; \D{2,1} = -1; \D{3,1} = 1; 
%   \E{1,1} = -2; \E{2,1} = 1; \E{3,1} = 1; 
%   \F{1,1} = 2; \F{2,1} = 1; \F{3,1} = 1; 
% }
% \tikzset{
%   conics/define/five points={\A,\B,\C,\D,\E,\Ca},
%   conics/define/five points={\A,\B,\C,\D,\F,\Cb},
%   conics/intersect cc={\Ca,\Cb,\Cdeg,\l,\m},
% }

\tikzmath{
  % 2x^2+xy+y^2-4x-4y=0
  \coefa{1} = 2; \coefa{2} = 2; \coefa{3} = 1; 
  \coefa{4} = -4; \coefa{5} = -4; \coefa{6} = 0;  
  % 4x^2+xy+y^2-8x-4y=0
  \coefb{1} = 4; \coefb{2} = 1; \coefb{3} = 1; 
  \coefb{4} = -8; \coefb{5} =-4; \coefb{6} = 0; 
}

\tikzset{
  conics/define/quadratic={\coefa,\Ca},
  conics/define/quadratic={\coefb,\Cb},
  conics/intersect cc={\Ca,\Cb,\Cdeg,\l,\m},
  mat/show={\Ca,3,3},
  mat/show={\Cb,3,3},
  mat/show={\Cdeg,3,3},
  mat/show={\l,3,1},
  mat/show={\m,3,1},
}

\begin{tikzpicture}
  \clip (-6,-6) rectangle (6,6);% 裁剪区域
  \axes[grid] (-5:5,-5:5);
  \conic[draw=navy] (\Ca);
  \conic[draw=teal] (\Cb);
  \segment[draw=red,clip=(-5:5,-5:5)](\l);
  \segment[draw=red,clip=(-5:5,-5:5)](\m);

  \coordinate (P1) at (0,0);
  \coordinate (P2) at (2,0);
  \coordinate (P3) at (2,2);
  \coordinate (P4) at (0,4);

  \foreach \x in {30,60,...,90}{
    \tikzset{%color
      mat/combine={\Ca,\Cb,3,3,\x/100,-0.75,\C},
    }
    \conic[draw=teal!\x] (\C);
  }

  \foreach \x in {30,60,...,90}{
    \tikzset{%color
      mat/combine={\Ca,\Cb,3,3,\x/100,0.75,\C},
    }
    \conic[draw=red!\x] (\C);
  }

  \foreach \p/\placement in {P1/above, P2/above, P3/above, P4/above}{
    \fill[red] (\p) circle (1.5pt);
    \draw (\p) node[\placement] {$\p$};
  }
\end{tikzpicture}

\end{document}  
